\section*{Data and code availability}
\addcontentsline{toc}{section}{Data and code availability}

All code, the generator, the evidence files and this paper's sources are in the repository at
\url{https://github.com/mariusbayizere/fraudshield}, on the branch \texttt{m11/paper}; the default
branch, \texttt{main}, does not yet hold the paper. The repository is licensed Apache-2.0; the dataset release terms are
CC BY 4.0 (\texttt{dataset/release/}). The dataset is not distributed as a file: seed \nSeed{} and the
generator at the milestone tag \texttt{m4-complete} regenerate it, as a regeneration at that tag has
shown by matching the fingerprint, which is evidence of reproduction rather than a comparison of every
row (Section~\ref{sec:repro}), and its fingerprint (Table~\ref{tab:glance}) identifies
the draw. Every number in the paper's text is mapped to its recorded source in \texttt{numbers.tex} and
checked against it by \texttt{audit\_numbers.py}, whose output is \texttt{number\_audit.md}; every
reference is checked in \texttt{reference\_audit.md}; the figures are regenerated from the audited
numbers by \texttt{figures/make\_figures.py}. No individual-level data about any real customer is used anywhere;
the \nProvSourced{} sourced generator parameters come from published aggregate statistics (central-bank
payment reports, a census and a financial-inclusion survey), from reference lists (ISO 4217, the IANA
time zone database, World Bank exchange rates and regional-bloc membership pages) and from one
school's calendar (Section~\ref{sec:dataset}).
